% Part: computability
% Chapter: computability-theory
% Section: applications-fixed-point

\documentclass[../../../include/open-logic-section]{subfiles}

\begin{document}

\olfileid[pt]{cmp}{thy}{apf}
\olsection{Aplicação do teorema do ponto fixo}

O teorema do ponto fixo essencialmente nos permite definir funções
parcialmente computáveis em termos de seus índices. Por exemplo,
podemos encontrar um índice $e$ tal que, para todo $y$,
\[
\cfind{e}(y) = e + y.
\]
Como outro exemplo, pode-se usar a demonstração do teorema do ponto
fixo para projetar um programa em Java ou C++ que imprima a si próprio.

Lembre-se de que, se, para cada $e$, tomarmos $W_e$ como o domínio de
$\cfind{e}$, então a sequência $W_0$, $W_1$, $W_2$,~\dots enumera os
conjuntos computavelmente enumeráveis. Alguns desses conjuntos são
computáveis. Pode-se perguntar se existe um algoritmo que receba como
entrada um valor $x$ e, se $W_x$ for computável, devolva um índice de
sua função característica. A resposta é ``não'': não existe tal
algoritmo.

\begin{thm}
Não existe função parcialmente computável $f$ com a seguinte
propriedade: sempre que $W_e$ é computável, $f(e)$ está definida e
$\cfind{f(e)}$ é sua função característica.
\end{thm}

\begin{proof}
Seja $f$ uma função computável qualquer; construiremos um $e$
tal que $W_e$ é computável, mas $\cfind{f(e)}$ não é sua função
característica. Usando o teorema do ponto fixo, podemos encontrar um
índice $e$ tal que
\[
\cfind{e}(y) \simeq
\begin{cases}
0 & \text{se $y=0$ e $\cfind{f(e)}(0) \fdefined = 0$} \\
\text{indefinida} & \text{caso contrário.}
\end{cases}
\]
Isto é, $e$ é obtido aplicando o teorema do ponto fixo à função
definida por
\[
g(x,y) \simeq
\begin{cases}
0 & \text{se $y=0$ e $\cfind{f(x)}(0) \fdefined = 0$} \\
\text{indefinida} & \text{caso contrário.}
\end{cases}
\]
Informalmente, podemos ver que $g$ é parcialmente computável da
seguinte maneira: nas entradas $x$ e $y$, o algoritmo verifica primeiro
se $y$ é igual a~$0$. Se for, o algoritmo calcula $f(x)$ e então usa
a máquina universal para calcular $\cfind{f(x)}(0)$. Se esta última
computação parar e devolver~$0$, o algoritmo devolverá~$0$; caso
contrário, o algoritmo não parará.

Observe agora que, se $\cfind{f(e)}(0)$ está definida e é igual a $0$,
então $\cfind{e}(y)$ está definida exatamente quando $y$ é igual a $0$,
de modo que $W_e = \{ 0 \}$. Se $\cfind{f(e)}(0)$ não estiver definida,
ou estiver definida mas não for igual a $0$, então
$W_e = \emptyset$. Em qualquer dos casos, $\cfind{f(e)}$ não é a
função característica de~$W_e$, pois dá a resposta errada na entrada
$0$.
\end{proof}

\end{document}

